% Generated draft: requires mathematical and source review.
\documentclass[11pt]{article}
\usepackage[margin=25mm]{geometry}
\usepackage{fontspec}
\setmainfont{Latin Modern Roman}
\usepackage{amsmath,amssymb,mathtools,unicode-math}
\setmathfont{Latin Modern Math}
\usepackage{enumitem,needspace,xurl,hyperref}
\setlength{\parindent}{0pt}
\setlength{\parskip}{5pt}
\setlength{\emergencystretch}{3em}
\newcommand{\N}{\mathbb{N}}
\newcommand{\Z}{\mathbb{Z}}
\newcommand{\Q}{\mathbb{Q}}
\newcommand{\R}{\mathbb{R}}
\newcommand{\C}{\mathbb{C}}
\newcommand{\F}{\mathbb{F}}
\newcommand{\A}{\mathbb{A}}
\newcommand{\PP}{\mathbb{P}}
\newcommand{\E}{\mathbb{E}}
\newcommand{\eps}{\varepsilon}
\newcommand{\dd}{\,\mathrm{d}}
\begin{document}
\section*{UnsolvedMath ID: 800919}
\textbf{Source problem code:} catalogue entry
\subsection{Definitions and mathematical statement}
For a positive integer \(n\), the supplied question concerns the moduli space denoted by \(M^{K3}_{2n}\). Under the conventional interpretation as the moduli space of complex primitively polarized \(K3\) surfaces, its points represent pairs \((S,H)\), where \(S\) is a smooth projective complex \(K3\) surface and \(H\) is an ample line bundle with primitive first Chern class and intersection number \(H^2=2n\). The supplied record does not specify the ground field, primitivity convention, stack versus coarse moduli space, or the quantification over \(n\); this interpretation is therefore provisional.

Under this interpretation, a Noether--Lefschetz divisor is an irreducible divisorial component of a locus where \(\operatorname{Pic}(S)\) contains a class \(\beta\) independent of \(H\) over \(\Q\), with prescribed intersection numbers \(H\cdot\beta\) and \(\beta^2\). Let \(\mathcal D_{\mathrm{NL}}\) denote the collection of these divisors. Here \(\operatorname{Pic}(M^{K3}_{2n})\) is the group of line bundles on the chosen moduli object. Interpreting each divisor as a rational Picard class, the question is whether
\[
\operatorname{span}_{\Q}\bigl\{[D]:D\in\mathcal D_{\mathrm{NL}}\bigr\}
=
\operatorname{Pic}(M^{K3}_{2n})\otimes_{\Z}\Q.
\]
Equivalently, must every rational Picard class be a finite rational linear combination of Noether--Lefschetz divisor classes? The choice of moduli object and the interpretation of divisor classes must be fixed for a definitive formulation.

The supplied motivation considers a compact Calabi--Yau threefold \(X\), a \(K3\) fibration \(\pi:X\to\PP^1\), and an ample line bundle \(L\) on \(X\). It states that this family determines a morphism from \(\PP^1\) to the moduli of polarized \(K3\) surfaces. The record gives no hypotheses on singular fibers or treatment of degenerations, so the asserted morphism over the entire base needs qualification; a smooth polarized family gives the corresponding moduli map over its smooth locus.

As historical background, the supplied source attributes to [mpnl] a calculation of the Gromov--Witten theory of \(X\) in \(\pi\)-fiber classes in terms of the Noether--Lefschetz numbers of \(\pi\) and the Katz--Klemm--Vafa conjecture [kkv] concerning \(\lambda_g\) integrals in the reduced Gromov--Witten theory of a fixed \(K3\) surface. It reports that the conjecture is proved for all classes in genus zero in [kmps] and for primitive classes in all genera in [mpt]. These are supplied historical source statements, not independently verified literature assessments.
\subsection{Short English statement}
Do Noether--Lefschetz divisors generate every rational Picard class on the moduli space \(M^{K3}_{2n}\) of polarized \(K3\) surfaces? The precise moduli convention and the intended quantification over \(n\) require clarification. The registry labels the problem open; that label has not been independently verified.
\subsection{Sources}
\begin{itemize}
\item Local UnsolvedMath registry, supplied complete record: permanent ID \(800919\), exact source problem code \texttt{catalogue entry}, title ``AGNES Open Problems in Algebraic Geometry --- Q7.'' Attribution: Pandharipande; proposed year: 2010. The supplied registry status is \texttt{open}, not an independently verified mathematical or literature assessment. The record was created and updated on 2026-10-16.
\item Supplied source list: Pandharipande, \emph{Open Problems in Enumerative Algebraic Geometry (AGNES)} (2010), source item Q7 (\#6), \url{https://arxiv.org/abs/1004.3259}. The record reports access on 2026-10-16 and extraction by \texttt{arxiv\_latex}; these provenance claims are supplied and have not been independently verified. Upstream revision: \texttt{d371ee917ac9647fbb06cd6fc529deec6f3dda82}.
\item References appearing in the supplied background: [mpnl], on fiber-class Gromov--Witten theory and Noether--Lefschetz numbers; [kkv], the Katz--Klemm--Vafa conjecture; [kmps], the reported genus-zero result for all classes; and [mpt], the reported all-genus result for primitive classes. Full bibliographic entries and individual links for these labels were not supplied, so their identities and precise hypotheses remain unresolved.
\end{itemize}
\end{document}
